\section{Justificación}

\subsection{Justificación técnica}

Desde la perspectiva técnica, el proyecto propone organizar la información de los servicios de soporte mediante un modelo estructurado de datos y la definición clara de estados de trabajo. Asimismo, contempla el diseño de un asistente conversacional basado en inteligencia artificial para apoyar la atención de consultas frecuentes sobre el estado de los equipos. Esta herramienta actuará como un canal complementario de información, sin sustituir la revisión técnica ni el diagnóstico especializado a cargo del personal calificado.

\subsection{Justificación social}

En el ámbito social, la propuesta busca facilitar el acceso de los clientes a la información sobre sus servicios de soporte técnico, reduciendo la incertidumbre durante el periodo de reparación. El diseño del sistema podría favorecer una comunicación más clara entre el usuario y el prestador del servicio, promoviendo acuerdos transparentes sobre los trabajos requeridos.

\subsection{Justificación económica}

Desde el punto de vista económico, el diseño propuesto busca respaldar la gestión organizativa de los prestadores de servicio. La centralización de la información y la automatización de respuestas a consultas frecuentes podrían contribuir a reducir el tiempo dedicado a la atención de llamadas o mensajes repetitivos, permitiendo un uso más eficiente del tiempo de trabajo.

\subsection{Justificación medioambiental}

En el aspecto ambiental, la digitalización de los registros de recepción y del seguimiento de las órdenes de trabajo podría reducir progresivamente la necesidad de utilizar comprobantes en papel. Adicionalmente, el adecuado mantenimiento de los equipos informáticos constituye una práctica que contribuye a prolongar su vida útil, lo que representa un beneficio indirecto potencial en favor del aprovechamiento de los recursos tecnológicos.
